\chapter{Conclusion}
\label{ch:conclusion}

\section{What was established}

This thesis set out to test whether per-layer local optima compose in
multi-tenant hybrid serving, and answered with a built system and a
pre-registered measurement campaign. The answer: they do not compose, and
the gap is large. A controller that plans replicas, semantic-cache budget,
and model tier together beats every tuned per-layer baseline on the
composite objective in every campaign that tested it, is the only
Pareto-undominated system in a 1{,}800-run matrix, and on replayed real
demand delivers $-70\%$ cost against tuned reactive scalers with a
disclosed $+0.07$ violation trade. Removing jointness alone costs
$+2884\%$ on the cost term --- the composition gap made into a number.

Around that central result: an honest SLO account (parity at
$-43\ldots{-48\%}$ cost, a confirmed forecasting mechanism, two published
confirmatory failures, and a stopping rule); real-data cache results with a
quality denominator no surveyed cache paper publishes; fairness wins over
two dedicated baselines at the stated altitude with the controller's own
fairness term honestly nulled; a security frontier showing per-tenant cache
partitioning dominates padding and jitter defenses against a demonstrated
AUC-0.88 timing side channel; two construction-level guarantees on the
deployed solver; and a working Kubernetes-native implementation held to CI
gates throughout.

The methodological result may outlast the numbers: seven pre-registrations
pushed before their runs, immutable closed campaigns, effect sizes over
p-values, substrate separation, and a published honest-nulls ledger ---
demonstrated in a systems evaluation at thesis scale, in public git
history.

\section{Future work}

\begin{itemize}
  \item \textbf{A three-knob live substrate.} The completed live campaign
        (Section~\ref{sec:eval-live}) recorded ordinal disagreement with
        live parity between the arms --- because the live data plane
        exposes only the replica lever. Extending it with a real semantic
        cache and tier-differentiated service is the highest-priority
        follow-up: it is the experiment that could confirm, or refute, the
        joint-control ranking outside the simulator.
  \item \textbf{Token-level demand signals.} Swap the request-rate
        \texttt{DemandSource} for TTFT/queue-depth/KV-pressure signals from
        an llm-d-class engine, keeping the planner unchanged --- the
        designed migration path onto the 2026 substrate
        (Section~\ref{sec:bg-substrate}).
  \item \textbf{p99 live restatement.} The live path measures wall-time
        latency; extending the ordinal protocol to p99 removes the
        simulator's p95 scope limit.
  \item \textbf{Planner scaling past 128 tenants.} Planning cells,
        incremental fairness partial sums, and a faster inner loop are the
        named levers against the measured 1.45 growth exponent.
  \item \textbf{Adopting the empirical cache curve.} Feeding the measured
        $h(K)$ fit into the model constants is a new pre-registered
        experiment by the project's own rules --- deliberately left for one.
  \item \textbf{Prefix/KV-cache layers and token-level fairness.} Both sit
        below \polyforge{}'s altitude and compose with it; a joint
        portfolio-plus-engine study is the natural next thesis.
\end{itemize}

\section{Closing remark}

The strongest sentence this project can write is not a percentage. It is
that every percentage in it can be traced --- to a protocol committed before
the run, a script that regenerates it, and a git history that would have
exposed any other way of obtaining it. That standard is portable to any
systems evaluation, and it is the practice this thesis most hopes to see
reused.
